\section{L'Espace personnel et les services des chercheurs}
La refonte de la bibliothèque numérique de Cujas doit s'affranchir d'une consultation passive de l'interface en un véritable espace de travail personnalisé tout en y adossant des outils d'analyse scientifique avancées.

\subsection{L'intégratiion du standard IIIF et de la visionneuse Mirador}
Sur CujasNum, à l'origine, la consultation des ouvrages reposait sur une visionneuse JPEG, de bonnes qualités. De plus, elle offrait également la possibilité de télécharger le document consulté sans format PDF sans offrir différentes options de téléchargements ou d'autres formats. Un des chercheurs expliquait qu'il était peu courant qu'une bibliothèque offre différents options de téléchargement, qu'il trouvait cela dommage car par moment il aurait bien aimé pouvoir choisir uniquement de télécharger les pages intéressantes pour son étude. Une doctorante mentionne le fait que ce serait bien d'ajouter en plus des options de téléchargement \enquote{la possibilité de conserver la page de titre de l'ouvrage}. De cette façon, l'usager se souvient d'où ces pages ont été extraites.

Pour répondre à ces attentes, le prototype Codex, développé par la bibliothèque Cujas, sépare la couche de diffusion de celle de conservation, grâce à l'implémentation d'un serveur IIIF. Ce standard permet la diffusion dynamique d'image accompagnée d'outils scientifiques, proposant la possibilité de réaliser des captures d'écran, des annotations.

L'implémentation de la visionneuse Mirador offre l'opportunité aux chercheurs de la comparaison des documents au sein d'un même espace de travail. Les chercheurs peuvent désormais consulter plusieurs ouvrages afin de mener des analyses comparatives. Un autre des avantages est la possibilité d'utiliser un zoom profond sur une image exposée en haute définition. Elle offre aussi des réglages de rotations d'images, de contrastes et de colorimétrie.

\subsection{L'espace personnel du chercheur}
L'enquête de terrain et les entretiens menés auprès des usagers révèlent leur souhait d'avoir leurs propres espaces de  travail pour mener à bien leurs recherches. Sous CujasNum, il avait pris l'habitude de télécharger chaque document ne pouvant les conserver une fois le navigateur fermé.

Pour ancrer durablement l'usager dans la bibliothèque numérique de Cujas, l'implémentation d'un espace personnel est une priorité fonctionnelle majeure. Cet espace proposera des services simples et intuitifs comme la gestion des favoris, l'archivage des requêtes, et la recherche ciblée dans les sélections. L'espace personnel offre aussi la possibilité de rédiger des notes de travail ainsi que des annotations personnelles. Comme l'a expliqué Elina Leblanc, l'espace personnel répond à un besoin d'archivage d'usage. Il s'agit comme \enquote{une barrière à l'oubli} et génère \enquote{une économie de temps} cruciale pour les chercheurs.

\subsection{L'exposition des données}
Afin d'accompagner les usagers dans leurs rédactions, le prototype Codex doit offrir la possibilité des services d'exportation des données et d'analyse linguistiques aux standards du web sémantique.

Les chercheurs interrogés sont pour l'ajout de l'export de références bibliographiques normalisées, comme Zotero, leur permettant ainsi de gagner du temps. comme le souligne Antoine Courtin, \enquote{que les pratiques des chercheurs en termes de données, se limitent bien souvent à copier des données dans un document de traitemenet}. De plus, grâce à l'implémentation du module open-source \textit{Bibliography} d'Omeka S, la plateforme rendra possible l'exportation instantané de la notice bibliographique vers des gestionnaires de références, comme Zotero ou EndNote, sous des formats interopérables standardisés tels que BibTex, RIS ou JSON-LD.

L'intégration du format XML-TEI s'impose comme le standard pour encader la structure logique des textes juridiques, facilitant ainsi la reconnaissance des entités nommées (personnalités citées, lieux, date) et permet la fouille de texte. Comme l'expliquent Emmanuelle Bermès et Gautier Poupeau, la fouille de texte est rendue possible grâce à la disponibilité des corpus. \enquote{La \enquote{fouille de textes et de données}(en anglais \textit{text and data mining} ou TDM) embrasse un nouveau paradigme dans l'accès aux documents ou aux sources, qu'on pourrait désigner sous le nom de lecture distante, dans la lignée de Franco Moretti}.L'intérêt fondamental du TDM permet d'exploiter la masse de textes désormais numérique s'affranchissant des limites physiques de la lecture traditionnelle.

En dotant l'usager d'outil de pointe comme une visionneuse Mirador, un compte personnel et des formats structurés, la bibliothèque numérique devient un véritable outil pour mener à bien les projets en humanités numériques. Afin de mettre en avant cette collaboration entre les chercheurs et la bibliothèque, l'infrastructure doit alors s'accompagner d'une médiation intellectuelle. 